\documentclass[10pt,a4paper]{amsart}
\usepackage[margin=19mm]{geometry}
\usepackage{amsmath,amssymb,mathtools}
\usepackage[scr=rsfs,cal=euler]{mathalfa}
\usepackage{xfrac}
\usepackage[unicode,hidelinks,bookmarks=false]{hyperref}
\newcommand{\ie}{i.e.\ }
\newcommand{\cf}{cf.\ }
\newcommand{\resp}{respectively\ }
\newcommand{\Subsection}{\subsection}
\providecommand{\nobreakdash}{\nobreak\hbox{-}}
\setlength{\emergencystretch}{2em}
\multlinegap0pt
\usepackage{amssymb}
\usepackage{mathtools}
\usepackage{xcolor}
\usepackage{float}
\usepackage[shortlabels]{enumitem}
\newtheorem{lemma}{Lemma}
\newtheorem{proposition}{Proposition}
\newcommand{\abs}[1]{\left\lvert#1\right\rvert}
\newcommand{\const}[1]{%
  \csname const@#1\endcsname%
}
\newcommand{\constapp}[1]{%
  \csname constapp@#1\endcsname%
}
\newcommand{\eps}{\varepsilon}
\renewcommand{\geq}{\geqslant}
\renewcommand{\leq}{\leqslant}
\newcommand{\prooflabel}[2]{%
  \def\@currentlabelname{#2}%
  \label{#1}%
}
\title{Wildfire in a narrow gully: a geometric reduction approach}
\author{Lorenzo De Gaspari, Serena Dipierro, Enrico Valdinoci}
\date{Source: arXiv:2604.12192v1 (CC BY 4.0)}
\begin{document}
\maketitle

\noindent\emph{Source and licence.} Wildfire in a narrow gully: a geometric reduction approach. Version arXiv:2604.12192v1 (CC BY 4.0), distributed by arXiv under the Creative Commons Attribution 4.0 International licence, http://creativecommons.org/licenses/by/4.0/, as recorded in the arXiv licence metadata for this exact version and shown on the abstract page https://arxiv.org/abs/2604.12192v1. Full licence notice: \texttt{licenses/arxiv-2604.12192-CC-BY-4.0.txt}.\par\smallskip

\noindent\textit{Editorial scope.} Counted unit: the article's lemma LEM::PiecewiseLipschitzToGlobalLipschitz (source0.tex lines 1050--1065). The article's complete original proof of it is delivered with it (source0.tex lines 2626--2786, one \texttt{proof} environment of 161 lines, which the article places away from the statement). No mathematics is written, repaired or re-derived: the statement and the proof are the article's own lines, in the article's own notation, and no retained line is substituted or re-flowed.  Retained uncounted as the source-local context the proof needs: lines 976--978; lines 2129--2136; lines 2436--2441; lines 2438--2440.  Nothing was omitted from any retained span, so the package carries no omission record.  No retained line contains a citation, so the wrapper carries no bibliography and no bibliography entry.  No retained line contains a figure environment, an image-inclusion command or drawing code, so no image is delivered and the package declares no asset dependency.  Editorial formatting only, in addition: an AMS \texttt{amsart} preamble that restates the article's own declarations and macros used by the retained lines, namely the environments \texttt{lemma}, \texttt{proposition} on separate counters, together with the standard packages those lines need.  The retained environments print in this document as Lemma 1, Proposition 1 in order of appearance, whereas the article prints the same items with the numbering of its own full document.  The complete native source of the article as distributed in the arXiv e-print for this exact version is bundled unchanged as \texttt{sources/198393/source0.tex}.
    \begin{equation}\label{EQ::VkDefinition}
        V_m \coloneq \bigcup_{s \in [(-1+2m)\eps, (1+2m)\eps]} \mathcal{S}(s).
    \end{equation}

\begin{proposition}\label{PROP::HolderEmbeddingA}
    Let~$X$ be either a domain or time cylinder. Let~$a \in (0,3)$,~$a^\prime \in [0,a)$, and~$b \in [-a^{\prime},+\infty)$. %chktex 9
    
    Then, there exists a constant~$\constapp{CONST::HolderEmbeddingA}>0$, which depends only on~$X$,~$a$ and~$a^\prime$, such that, for every~$u \in \mathcal{H}_a^{(b)}(X)$,
    \begin{equation}\label{EQ::HolderEmbeddingA}
        \abs{u}_{a^\prime; X}^{(b)} \leq \constapp{CONST::HolderEmbeddingA} \abs{u}_{a; X}^{(b)}.
    \end{equation}
\end{proposition}

\begin{lemma}\label{LEM::Quasiconvexity}
    There exists a constant~$\constapp{CONST::Quasiconvexity} >0$, which depends only on~$n$,~$\mathcal{S}$ and~$L$, such that, for every~$x,x^\prime \in \Omega_L$, there exists a~$C^1$ curve~$\gamma \colon [0,1] \to \Omega_L$ such that
    \begin{equation}\label{EQ::QuasiconvexityInequality}
        \gamma(0)=x,\qquad\gamma(1)=x^\prime,\qquad\abs{x-x^\prime}\leq\operatorname{length}(\gamma)\leq \constapp{CONST::Quasiconvexity} \abs{x-x^\prime}.
    \end{equation}
\end{lemma}

    \begin{equation}
        \gamma(0)=x,\qquad\gamma(1)=x^\prime,\qquad\abs{x-x^\prime}\leq\operatorname{length}(\gamma)\leq \constapp{CONST::Quasiconvexity} \abs{x-x^\prime}.
    \end{equation}

\begin{lemma}\label{LEM::PiecewiseLipschitzToGlobalLipschitz}
    Let~$b \in \{-1\} \cup [0,+\infty)$ and~$u \in C(\Omega_\tau \cup \mathcal{N}_\tau)$ be such that%chktex 9
    \begin{equation}
        u \in \mathcal{H}_1^{(b)}(V_m)\qquad \text{for every }m \in \{-\mathfrak{K},\dots,\mathfrak{K}\}.
    \end{equation}

    Then, there exists a constant~$\const{CONST::PiecewiseLipschitzToGlobalLipschitz} > 0$, which depends only on~$n$,~$\mathcal{S}$,~$L$, and~$b$, such that, for every $a \in (0,1)$,
    \begin{equation}\label{EQ::PiecewiseLipschitzToGlobalHolder}
        \abs{u}_{a; \Omega_\tau}^{(\max\{-a,b\})} \leq \const{CONST::PiecewiseLipschitzToGlobalLipschitz} \max_{m \in \{-\mathfrak{K},\dots,\mathfrak{K}\}} \abs{u}_{1; V_m}^{(b)}.
    \end{equation}

    Moreover, if~$\nabla u \in C(\Omega_\tau \cup \mathcal{N}_\tau)$, then
    \begin{equation}\label{EQ::PiecewiseLipschitzToGlobalLipschitz}
        \abs{u}_{1; \Omega_\tau}^{(b)} \leq \const{CONST::PiecewiseLipschitzToGlobalLipschitz} \max_{m \in \{-\mathfrak{K},\dots,\mathfrak{K}\}} \abs{u}_{1; V_m}^{(b)}.
    \end{equation}
\end{lemma}

\begin{proof}[Proof of Lemma~\ref{LEM::PiecewiseLipschitzToGlobalLipschitz}]\prooflabel{PRF::PiecewiseLipschitzToGlobalLipschitz}{Proof of Lemma~\ref{LEM::PiecewiseLipschitzToGlobalLipschitz}}
    Let~$\delta > 0$, and~$x,x^\prime \in I_\delta(\Omega_\tau)$, with~$x \neq x^\prime$. Consider the following two cases: either~$\abs{x-x^\prime} < \frac{\delta}{2}$, or~$\abs{x-x^\prime} \geq \frac{\delta}{2}$.

    If~$\abs{x-x^\prime} < \frac{\delta}{2}$, then
    \begin{equation}\label{EQ::SegmentIsInOmegaDeltaLipschitz}
        E \coloneq \bigg\{(1-\mu)x + \mu x^\prime \colon \mu \in (0,1)\bigg\} \subset \overline{B}_{\frac{\delta}{2}}(x) \subset I_{\frac{\delta}{2}}(\Omega_\tau),
    \end{equation}
    therefore, the~$n$-th Fermi coordinate (i.e., the second component of~$\Phi^{-1}$) is well defined and continuous on~$E$.

    From~\eqref{EQ::VkDefinition}, we see that
    \begin{equation}\label{EQ::OmegaDeltaVkPartitionLipschitz}
        I_{\frac{\delta}{2}}(\Omega_\tau) = \bigcup_{m \in \{-\mathfrak{K}, \dots, \mathfrak{K}\}}I_{\frac{\delta}{2}}(V_m),
    \end{equation}
    therefore, there exist~$m_x$ and~$m_{x^\prime}$ such that~$x \in I_{\frac{\delta}{2}}(V_{m_x})$ and~$x^\prime \in I_{\frac{\delta}{2}}(V_{m_{x^\prime}})$. Without loss of generality, we assume that~$m_x \leq m_{x^\prime}$.

    If~$m_x = m_{x^\prime}$, then~$x,x^\prime \in I_{\frac{\delta}{2}}(V_{m_x})$, thus
    \begin{equation}\label{EQ::xxPrimeNearSameVkLipschitz}
        \abs{u(x)-u(x^\prime)} \leq \abs{u}_{1; I_{\frac{\delta}{2}}(V_{m_x})} \abs{x-x^\prime} \leq \max_{m \in \{-\mathfrak{K},\dots,\mathfrak{K}\}} \abs{u}_{1; I_{\frac{\delta}{2}}(V_m)}\abs{x-x^\prime}.
    \end{equation}

    If, instead,~$m_x < m_{x^\prime}$, we let
    \begin{equation}\label{OLD::EQ::KdefinitionAndBound}
        K \coloneq m_{x^\prime} - m_x \leq 2\mathfrak{K} \leq \frac{L}{\eps}
    \end{equation}
    and~${\{\mu_j\}}_{j=1}^K$ be such that
    \begin{equation}\label{EQ::mujSubsegmentDefinitionLipschitz}
        \mu_j \coloneq \inf\left\{\mu \in [0,1] \colon (1-\mu)x + \mu x^\prime \in I_{\frac{\delta}{2}}(V_{m_x + j})\right\}.
    \end{equation}
    Each~$\mu_j$ is well defined because of the continuity of the~$n$-th Fermi coordinate and because, thanks to~\eqref{EQ::SegmentIsInOmegaDeltaLipschitz}, the segment~$E$ is fully contained in~$I_{\frac{\delta}{2}}(\Omega_\tau)$. Hence, letting~$x_j \coloneq (1-\mu_j)x + \mu_j x^\prime$ for each~$j \in \{1,\dots,m_{x^\prime}-m_x\}$, we have
    \begin{equation}
        x_j \in I_{\frac{\delta}{2}}(V_{m_x + j - 1})\cap I_{\frac{\delta}{2}}(V_{m_x + j}),\qquad \text{for every }j \in \{1,\dots,m_{x^\prime}-m_x\}.
    \end{equation}
    Thus,
    \begin{equation}
        \begin{split}
            |u&(x)-u(x^\prime)| \leq \abs{u(x)-u(x_1)} + \sum_{j=1}^{K-1} \abs{u(x_j)-u(x_{j+1})} + \abs{u(x_K)-u(x^\prime)}\\
                &\leq \abs{u}_{1; I_{\frac{\delta}{2}}(V_{m_x})} \abs{x-x_1} + \sum_{j=1}^{K-1} \left( \abs{u}_{1; I_{\frac{\delta}{2}}(V_{m_x + j})} \abs{x_j-x_{j+1}}\right)+\abs{u}_{1; I_{\frac{\delta}{2}}(V_{m_{x^\prime}})} \abs{x_K - x^\prime}\\
                &\leq \max_{m \in \{-\mathfrak{K},\dots,\mathfrak{K}\}} \abs{u}_{1; I_{\frac{\delta}{2}}(V_m)} \left( \mu_1 + \sum_{j=1}^{K-1} {(\mu_{j+1}-\mu_j)} + {(1-\mu_K)}\right)\abs{x-x^\prime}.
        \end{split}
    \end{equation}
    Hence, using~\eqref{EQ::mujSubsegmentDefinitionLipschitz} we have that
    \begin{equation}\label{EQ::xxPrimeNearGeneralLipschitz}
            |u(x)-u(x^\prime)| \leq \max_{m \in \{-\mathfrak{K},\dots,\mathfrak{K}\}} \abs{u}_{1; I_{\frac{\delta}{2}}(V_m)}\abs{x-x^\prime}.
    \end{equation}
    Thanks to~\eqref{EQ::xxPrimeNearSameVkLipschitz}, we know that~\eqref{EQ::xxPrimeNearGeneralLipschitz} holds true whenever~$\abs{x-x^\prime} < \frac{\delta}{2}$.

    We now focus on the case~$\abs{x-x^\prime} \geq \frac{\delta}{2}$. We first observe that, due to~\eqref{EQ::OmegaDeltaVkPartitionLipschitz},
    \begin{equation}\label{EQ::PiecewiseSupToGlobalLipschitz}
        \abs{u}_{0;I_\delta(\Omega_\tau)} = \max_{m \in \{-\mathfrak{K},\dots,\mathfrak{K}\}}\abs{u}_{0;I_\delta(V_m)}.
    \end{equation}
    Then, if~$b \geq 0$, we compute
    \begin{align}\label{OLD::EQ::xxprimeFarPositiveB}
        \abs{u(x)-u(x^\prime)} &\leq \abs{u(x)}+\abs{u(x^\prime)} \leq 2\abs{u}_{0;I_\delta(\Omega_\tau)}\\
            &\leq 2\max_{m \in \{-\mathfrak{K},\dots,\mathfrak{K}\}}\abs{u}_{0;I_\delta(V_m)}{\left(\frac{2\abs{x-x^\prime}}{\delta}\right)}\\
            &\leq 4{\delta^{-1}}\max_{m \in \{-\mathfrak{K},\dots,\mathfrak{K}\}}\abs{u}_{0;I_\delta(V_m)}\abs{x-x^\prime}.
    \end{align}

    Combining this and~\eqref{EQ::xxPrimeNearGeneralLipschitz} we find
    \begin{equation}\label{EQ::PiecewiseToGlobalDeltaLipschitz}
        \abs{u(x)-u(x^\prime)} \leq \left(\max_{m \in \{-\mathfrak{K},\dots,\mathfrak{K}\}} \abs{u}_{1; I_{\frac{\delta}{2}}(V_m)} + 4{\delta^{-1}}\max_{m \in \{-\mathfrak{K},\dots,\mathfrak{K}\}}\abs{u}_{0;I_\delta(V_m)}\right) \abs{x-x^\prime}.
    \end{equation}
    We now let~$a \in (0,1)$ and observe that
    \begin{equation}\label{EQ::LipschitzToHolderDistance}
        \abs{x-x^\prime} \leq {\operatorname{diam}(\Omega_L)}^{1-a} \abs{x-x^\prime}^a \leq \max\{1,\operatorname{diam}(\Omega_L)\} \abs{x-x^\prime}^a,
    \end{equation}
    hence, using also~\eqref{EQ::PiecewiseSupToGlobalLipschitz} and~\eqref{EQ::PiecewiseToGlobalDeltaLipschitz} we have that
    \begin{equation}
        \abs{u}_{a; I_{\frac{\delta}{2}}(\Omega_\tau)}^{(b)} \leq \max\{1,\operatorname{diam}(\Omega_L)\} \left(\max_{m \in \{-\mathfrak{K},\dots,\mathfrak{K}\}} \abs{u}_{1; I_{\frac{\delta}{2}}(V_m)} + (1+4{\delta^{-1}})\max_{m \in \{-\mathfrak{K},\dots,\mathfrak{K}\}}\abs{u}_{0;I_\delta(V_m)}\right).
    \end{equation}

    We use this information to discover that
    \begin{equation}\label{EQ::LongComputationBis}
        \begin{split}
            \abs{u}_{a; \Omega_\tau}^{(b)} &= \sup_{\delta > 0}\delta^{1+b} \abs{u}_{1; I_\delta(\Omega_\tau)}\\
                &\leq \max\{1,\operatorname{diam}(\Omega_L)\} \sup_{\delta > 0}\delta^{1+b}\left(\max_{m \in \{-\mathfrak{K},\dots,\mathfrak{K}\}} \abs{u}_{1; I_{\frac{\delta}{2}}(V_m)} + (1+4{\delta^{-1}})\max_{m \in \{-\mathfrak{K},\dots,\mathfrak{K}\}}\abs{u}_{0;I_\delta(V_m)}\right)\\
                &= \max\{1,\operatorname{diam}(\Omega_L)\} \left(\max_{m \in \{-\mathfrak{K},\dots,\mathfrak{K}\}}\sup_{\delta > 0}\delta^{1+b}\abs{u}_{1; I_{\frac{\delta}{2}}(V_m)} +\max_{m \in \{-\mathfrak{K},\dots,\mathfrak{K}\}} \sup_{\delta > 0}\delta^{1+b}\abs{u}_{0;I_\delta(V_m)}\right.\\
                &\qquad \left.+ 4\max_{m \in \{-\mathfrak{K},\dots,\mathfrak{K}\}} \sup_{\delta > 0}\delta^{b}\abs{u}_{0;I_\delta(V_m)}\right).
        \end{split}
    \end{equation}
    Here, we explicitly use the fact that~$b \geq 0$, so that the space~$\mathcal{H}_0^{(b)}$ is nontrivial, and its norm is well defined, obtaining
    \begin{equation}\label{EQ::PiecewiseRegularityToGlobalLongComputationLipschitz}
        \begin{split}
                \abs{u}_{a; \Omega_\tau}^{(b)}&\leq \max\{1,\operatorname{diam}(\Omega_L)\} \left(2^{1+b}\max_{m \in \{-\mathfrak{K},\dots,\mathfrak{K}\}}\sup_{\delta > 0}\delta^{1+b}\abs{u}_{1; I_\delta(V_m)}\right.\\
                &\qquad \left.+ \operatorname{diam}(\Omega_L)\max_{m \in \{-\mathfrak{K},\dots,\mathfrak{K}\}} \sup_{\delta \in (0,\operatorname{diam}(\Omega_L))}\delta^{b}\abs{u}_{0;I_\delta(V_m)}+4\max_{m \in \{-\mathfrak{K},\dots,\mathfrak{K}\}} \abs{u}_{0;V_m}^{(b)}\right)\\
                &\leq \max\{1,\operatorname{diam}(\Omega_L)\} \left(2^{1+b} + \constapp{CONST::HolderEmbeddingA}\left(\operatorname{diam}(\Omega_L) + 4\right)\right)\max_{m \in \{-\mathfrak{K},\dots,\mathfrak{K}\}} \abs{u}_{1;V_m}^{(b)},
        \end{split}
    \end{equation}
    where in the last inequality we also used Proposition~\ref{PROP::HolderEmbeddingA} with the choice~$a \coloneq 1$,~$a^\prime \coloneq 0$, and~$b \coloneq b$.

    If~$b=-1$, we have that~$u \in \mathcal{H}_{1}^{(-1)}(V_m) = \mathcal{H}_{1}(V_m)$ for every~$m \in \{-\mathfrak{K},\dots,\mathfrak{K}\}$. Notably, this means that~$u$ is~$C^1$ up to the Dirichlet boundary in each~$V_m$. %chktex 9

    We let~$\gamma$ be the curve joining~$x$ and~$x^\prime$ given by Lemma~\ref{LEM::Quasiconvexity} (namely, we have that~$\gamma(0) = x$ and~$\gamma(1) = x^\prime$). We assume that~$m_x < m_x^\prime$, otherwise~\eqref{EQ::xxPrimeNearSameVkLipschitz} holds. Then, we let
    \begin{equation}
        \sigma_j \coloneq \inf\left\{ \sigma \in [0,1] \colon \gamma(\sigma) \in V_{m_x+j} \right\},
    \end{equation}
    for every~$j \in \{1,\dots,m_{x^\prime}-m_x\}$. As with~\eqref{EQ::mujSubsegmentDefinitionLipschitz}, the values~$\sigma_j$ are well defined and such that
    \begin{equation}
        y_j \coloneq \gamma(\sigma_j)\in I_{\frac{\delta}{2}}(V_{m_x + j - 1})\cap I_{\frac{\delta}{2}}(V_{m_x + j}),\qquad \text{for every }j \in \{1,\dots,m_{x^\prime}-m_x\}.
    \end{equation}
    Thus,
    \begin{equation}
        \begin{split}
        |u(&x)-u(x^\prime)|\leq \abs{u(x)-u(y_1)} + \sum_{j=1}^{K-1} \abs{u(y_j)-u(y_{j+1})} + \abs{u(y_K)-u(x^\prime)}\\
            &\leq \abs{u}_{1; V_{m_x}} \abs{x-y_1} + \sum_{j=1}^{K-1} \left( \abs{u}_{1; V_{m_x+j}} \abs{y_j-y_{j+1}}\right) +\abs{u}_{1; V_{m_{x^\prime}}} \abs{y_K - x^\prime}\\
            &\leq \max_{m \in \{-\mathfrak{K},\dots,\mathfrak{K}\}} \abs{u}_{1; V_{m}} \left( {\left(\int_0^{\sigma_1}\abs{\dot{\gamma}(t)}\,dt\right)}+ \sum_{j=1}^{K-1} {\left(\int_{\sigma_j}^{\sigma_{j+1}}\abs{\dot{\gamma}(t)}\,dt\right)} + {\left(\int_{\sigma_K}^{1}\abs{\dot{\gamma}(t)}\,dt\right)}\right)\\
            & = \operatorname{length}(\gamma) \max_{m \in \{-\mathfrak{K},\dots,\mathfrak{K}\}} \abs{u}_{1; V_{m}}.
        \end{split}
    \end{equation}
    Therefore, thanks to~\eqref{EQ::QuasiconvexityInequality} and~$\abs{x-x^\prime} > \frac{\delta}{2}$, we have that
    \begin{equation}\label{OLD::EQ::xxprimeFarNegativeB}
        |u(x)-u(x^\prime)| \leq \constapp{CONST::Quasiconvexity} \max_{m \in \{-\mathfrak{K},\dots,\mathfrak{K}\}}\abs{u}_{1; V_{m}}\abs{x-x^\prime},
    \end{equation}
    which, together with~\eqref{EQ::xxPrimeNearGeneralLipschitz}, gives that
    \begin{equation}\label{EQ::LipschitzEstimateNegativeB}
        |u(x)-u(x^\prime)| \leq \left(\max_{m \in \{-\mathfrak{K},\dots,\mathfrak{K}\}} \abs{u}_{1; I_{\frac{\delta}{2}}(V_m)} + \constapp{CONST::Quasiconvexity} \max_{m \in \{-\mathfrak{K},\dots,\mathfrak{K}\}}\abs{u}_{1; V_{m}}\right)\abs{x-x^\prime}.
    \end{equation}
    We combine this new information with~\eqref{EQ::LipschitzToHolderDistance} to obtain
    \begin{equation}
        |u(x)-u(x^\prime)| \leq \max\{1,\operatorname{diam}(\Omega_L)\}\left(\max_{m \in \{-\mathfrak{K},\dots,\mathfrak{K}\}} \abs{u}_{1; I_{\frac{\delta}{2}}(V_m)} + \constapp{CONST::Quasiconvexity} \max_{m \in \{-\mathfrak{K},\dots,\mathfrak{K}\}}\abs{u}_{1; V_{m}}\right)\abs{x-x^\prime}^a.
    \end{equation}
    This and~\eqref{EQ::PiecewiseSupToGlobalLipschitz} show that
    \begin{equation}\begin{split}
            \abs{u}_{a; I_\delta(\Omega_\tau)} \leq &\max\{1,\operatorname{diam}(\Omega_L)\}\max_{m \in \{-\mathfrak{K},\dots,\mathfrak{K}\}} \abs{u}_{1; I_{\frac{\delta}{2}}(V_m)} + \max_{m \in \{-\mathfrak{K},\dots,\mathfrak{K}\}}\abs{u}_{0;I_\delta(V_m)}\\
                &+ \constapp{CONST::Quasiconvexity} \max\{1,\operatorname{diam}(\Omega_L)\}\max_{m \in \{-\mathfrak{K},\dots,\mathfrak{K}\}}\abs{u}_{1; V_{m}}.
        \end{split}
    \end{equation}
    Finally,
    \begin{equation}\label{EQ::HolderStitchingBetaLipschitz}
        \begin{split}
            \abs{u}_{a; \Omega_\tau} &= \sup_{\delta > 0} \abs{u}_{a; I_\delta(\Omega_\tau)}\\
                &\leq \sup_{\delta > 0}\left(\max\{1,\operatorname{diam}(\Omega_L)\}\max_{m \in \{-\mathfrak{K},\dots,\mathfrak{K}\}} \abs{u}_{1; I_{\frac{\delta}{2}}(V_m)} + \max_{m \in \{-\mathfrak{K},\dots,\mathfrak{K}\}}\abs{u}_{0;I_\delta(V_m)}\right.\\
                &+ \left.\constapp{CONST::Quasiconvexity} \max\{1,\operatorname{diam}(\Omega_L)\}\max_{m \in \{-\mathfrak{K},\dots,\mathfrak{K}\}}\abs{u}_{1; V_{m}}\right)\\
                &= \left(1+(1+\constapp{CONST::Quasiconvexity}) \max\{1,\operatorname{diam}(\Omega_L)\}\right)\max_{m \in \{-\mathfrak{K},\dots,\mathfrak{K}\}}\abs{u}_{1; V_{m}}.
        \end{split}
    \end{equation}

    We observe that if~$b \geq 0$, then~$\max\{-a,b\} = b$, hence~\eqref{EQ::PiecewiseRegularityToGlobalLongComputationLipschitz} gives~\eqref{EQ::PiecewiseLipschitzToGlobalHolder} with
    \[\const{CONST::PiecewiseLipschitzToGlobalLipschitz} \coloneq \max\{1,\operatorname{diam}(\Omega_L)\} \left(2^{1+b} + \constapp{CONST::HolderEmbeddingA}\left(\operatorname{diam}(\Omega_L) + 4\right)\right).\]
    If~$b = -1$, instead, we have that~$\max\{-a,b\} = -a$, and~\eqref{EQ::PiecewiseLipschitzToGlobalHolder} follows from~\eqref{EQ::HolderStitchingBetaLipschitz} upon choosing
    \[\const{CONST::PiecewiseLipschitzToGlobalLipschitz} \coloneq \left(1+(1+\constapp{CONST::Quasiconvexity}) \max\{1,\operatorname{diam}(\Omega_L)\}\right).\]

    It remains to prove~\eqref{EQ::PiecewiseLipschitzToGlobalLipschitz}. If~$b \geq 0$, from~\eqref{EQ::PiecewiseToGlobalDeltaLipschitz} we have that~$u$ is Lipschitz continuous in~$I_\delta(\Omega_\tau)$, with Lipschitz constant
    \[C_{\delta} \coloneq \left(\max_{m \in \{-\mathfrak{K},\dots,\mathfrak{K}\}} \abs{u}_{1; I_{\frac{\delta}{2}}(V_m)} + 4{\delta^{-1}}\max_{m \in \{-\mathfrak{K},\dots,\mathfrak{K}\}}\abs{u}_{0;I_\delta(V_m)}\right).\]
    Instead, if~$b = -1$, the Lipschitz continuity of~$u$ still holds in~$I_\delta(\Omega_\tau)$ as a consequence of~\eqref{EQ::LipschitzEstimateNegativeB} with
    \[C_{\delta} \coloneq \left(\max_{m \in \{-\mathfrak{K},\dots,\mathfrak{K}\}} \abs{u}_{1; I_{\frac{\delta}{2}}(V_m)} + \constapp{CONST::Quasiconvexity} \max_{m \in \{-\mathfrak{K},\dots,\mathfrak{K}\}}\abs{u}_{1; V_{m}}\right).\]

    Thanks to the assumption that~$\nabla u$ is continuous, we have that
    \begin{equation}
        \abs{\nabla u}_{0; I_\delta(\Omega_\tau)} \leq C_{\delta},
    \end{equation}
    therefore,
    \begin{equation}
        \abs{u}_{1; I_\delta(\Omega_\tau)} \leq \abs{u}_{0; I_\delta(\Omega_\tau)} + C_\delta
    \end{equation}
    and
    \begin{equation}
        \abs{u}_{1; \Omega_\tau}^{(b)} \leq \sup_{\delta > 0} \delta^{1+b} (\abs{u}_{0; I_\delta(\Omega_\tau)}+C_\delta).
    \end{equation}

    The proof is concluded upon noticing that, up to replacing~$\max\{1,\operatorname{diam}(\Omega_L)\}$ with~$1$, the computation is the same as~\eqref{EQ::LongComputationBis} and~\eqref{EQ::PiecewiseRegularityToGlobalLongComputationLipschitz} if~$b \geq 0$, and~\eqref{EQ::HolderStitchingBetaLipschitz} if~$b = -1$.
\end{proof}

\end{document}
